\chapter{作用量、正则结构与正规模}

\section{从 Newton 到 D'Alembert}
对 $N$ 个质点
\[
m_\alpha\ddot{\vb r}_\alpha=\vb F_\alpha+\vb R_\alpha,
\]
$\vb R_\alpha$ 为约束力。完整约束写成
$f_s(\vb r_1,\ldots,t)=0$。同一时刻满足线性化约束的虚位移满足
\[
\sum_\alpha\nabla_\alpha f_s\cdot\delta\vb r_\alpha=0.
\]
理想约束定义为
$\sum_\alpha\vb R_\alpha\cdot\delta\vb r_\alpha=0$。将 Newton 方程移项、点乘虚位移并求和：
\[
\sum_\alpha(\vb F_\alpha-m_\alpha\ddot{\vb r}_\alpha)
\cdot\delta\vb r_\alpha
=-\sum_\alpha\vb R_\alpha\cdot\delta\vb r_\alpha=0.
\]
这就是 D'Alembert 原理。它不是新的力学定律，而是 Newton 方程在允许虚位移子空间中的投影。

\section{广义坐标推导 Lagrange 方程}
取独立坐标 $q_i$：
\[
\vb r_\alpha=\vb r_\alpha(q_1,\ldots,q_s,t),\qquad
\delta\vb r_\alpha=\sum_i\pdv{\vb r_\alpha}{q_i}\delta q_i.
\]
定义广义力
\[
Q_i=\sum_\alpha\vb F_\alpha\cdot
\pdv{\vb r_\alpha}{q_i}.
\]
惯性项使用链式法则：
\[
\pdv T{\dot q_i}
=\sum_\alpha m_\alpha\dot{\vb r}_\alpha\cdot
\pdv{\dot{\vb r}_\alpha}{\dot q_i}
=\sum_\alpha m_\alpha\dot{\vb r}_\alpha\cdot
\pdv{\vb r_\alpha}{q_i}.
\]
对时间求导：
\[
\dv{}{t}\pdv T{\dot q_i}
=\sum_\alpha m_\alpha\ddot{\vb r}_\alpha\cdot
\pdv{\vb r_\alpha}{q_i}
+\sum_\alpha m_\alpha\dot{\vb r}_\alpha\cdot
\dv{}{t}\pdv{\vb r_\alpha}{q_i}.
\]
另一方面
\[
\pdv T{q_i}
=\sum_\alpha m_\alpha\dot{\vb r}_\alpha\cdot
\pdv{\dot{\vb r}_\alpha}{q_i}
=\sum_\alpha m_\alpha\dot{\vb r}_\alpha\cdot
\dv{}{t}\pdv{\vb r_\alpha}{q_i}.
\]
相减得关键恒等式
\[
\sum_\alpha m_\alpha\ddot{\vb r}_\alpha\cdot
\pdv{\vb r_\alpha}{q_i}
=\dv{}{t}\pdv T{\dot q_i}-\pdv T{q_i}.
\]
代回 D'Alembert 原理，独立 $\delta q_i$ 的系数必须为零：
\[
\dv{}{t}\pdv T{\dot q_i}-\pdv T{q_i}=Q_i.
\]
若 $Q_i=-\partial V/\partial q_i$，定义 $L=T-V$：
\[
\boxed{\dv{}{t}\pdv L{\dot q_i}-\pdv L{q_i}=0}.
\]

\section{作用量变分与端点条件}
\[
S[q]=\int_{t_1}^{t_2}L(q,\dot q,t)\dd t.
\]
取 $q_i\to q_i+\epsilon\eta_i$，固定端点
$\eta_i(t_1)=\eta_i(t_2)=0$：
\[
\delta S=\int\sum_i\left(
\pdv L{q_i}\eta_i+\pdv L{\dot q_i}\dot\eta_i\right)\dd t.
\]
第二项分部积分：
\[
\int\pdv L{\dot q_i}\dot\eta_i\dd t
=\left.\pdv L{\dot q_i}\eta_i\right|_{t_1}^{t_2}
-\int\dv{}{t}\pdv L{\dot q_i}\eta_i\dd t.
\]
边界项因固定端点为零，
\[
\delta S=\int\sum_i\left[
\pdv L{q_i}-\dv{}{t}\pdv L{\dot q_i}
\right]\eta_i\dd t.
\]
任意内部 $\eta_i$ 推出 Euler--Lagrange 方程。作用量通常只是驻值，不保证最小值。

\section{Noether 定理的有限自由度推导}
无穷小变换
$q_i\to q_i+\epsilon X_i(q,t)$ 使
\[
\delta L=\epsilon\dv F t.
\]
沿满足运动方程的轨道：
\[
\delta L=\sum_i\left(
\pdv L{q_i}X_i+\pdv L{\dot q_i}\dv{X_i}{t}\right)\epsilon.
\]
用
$\partial L/\partial q_i=\dot p_i$：
\[
\delta L=\epsilon\sum_i(\dot p_iX_i+p_i\dot X_i)
=\epsilon\dv{}{t}\sum_ip_iX_i.
\]
与 $\epsilon\dot F$ 比较：
\[
\boxed{\dv{}{t}\left(\sum_ip_iX_i-F\right)=0}.
\]
空间平移、旋转和时间平移分别给线动量、角动量和能量。

\section{Legendre 变换到 Hamilton 方程}
若 Hessian
$W_{ij}=\partial^2L/\partial\dot q_i\partial\dot q_j$ 可逆，
$p_i=\partial L/\partial\dot q_i$ 可反解 $\dot q_i(q,p,t)$。定义
\[
H(q,p,t)=\sum_ip_i\dot q_i-L.
\]
全微分：
\[
\dd H=\sum_i(\dot q_i\dd p_i+p_i\dd\dot q_i)
-\sum_i\left(\pdv L{q_i}\dd q_i+
\pdv L{\dot q_i}\dd\dot q_i\right)-\pdv L t\dd t.
\]
$p_i=\partial L/\partial\dot q_i$ 使 $\dd\dot q_i$ 项抵消；Lagrange 方程使
$\partial L/\partial q_i=\dot p_i$：
\[
\dd H=\sum_i(\dot q_i\dd p_i-\dot p_i\dd q_i)
-\pdv L t\dd t.
\]
与
$\dd H=\sum_i(\partial_qH\dd q_i+\partial_pH\dd p_i)+\partial_tH\dd t$
比较：
\[
\boxed{\dot q_i=\pdv H{p_i},\quad
\dot p_i=-\pdv H{q_i},\quad
\pdv Ht=-\pdv Lt}.
\]

\section{Poisson 括号与 Liouville 定理}
\[
\{f,g\}=\sum_i\left(
\pdv f{q_i}\pdv g{p_i}
-\pdv f{p_i}\pdv g{q_i}\right).
\]
链式法则配合 Hamilton 方程：
\[
\dot f=\sum_i(\partial_{q_i}f\,\dot q_i+
\partial_{p_i}f\,\dot p_i)+\partial_tf
=\{f,H\}+\partial_tf.
\]
相空间速度散度：
\[
\sum_i\left(
\pdv{\dot q_i}{q_i}+\pdv{\dot p_i}{p_i}\right)
=\sum_i\left(
\frac{\partial^2H}{\partial q_i\partial p_i}
-\frac{\partial^2H}{\partial p_i\partial q_i}\right)=0.
\]
因此相体积不变。概率密度满足
$\partial_t\rho+\{\rho,H\}=0$，是经典 Liouville 方程。

\section{正则变换和生成函数}
要求一阶作用量只差全导数：
\[
\sum_ip_i\dd q_i-H\dd t
=\sum_iP_i\dd Q_i-K\dd t+\dd F.
\]
若选择第二类生成函数 $F_2(q,P,t)$，
\[
\dd F_2=\sum_i\pdv{F_2}{q_i}\dd q_i+
\sum_i\pdv{F_2}{P_i}\dd P_i+\pdv{F_2}t\dd t.
\]
比较独立微分得到
\[
\boxed{p_i=\pdv{F_2}{q_i},\qquad
Q_i=\pdv{F_2}{P_i},\qquad
K=H+\pdv{F_2}t}.
\]
这比死记四类生成函数可靠：先写正则一形式，再按所选独立变量作 Legendre 变换。

\section{Hamilton--Jacobi 方程}
寻找正则变换使新 Hamiltonian $K=0$。令生成函数
$S(q,\alpha,t)$：
\[
p_i=\pdv S{q_i},\qquad
\boxed{\pdv St+H\left(q,\pdv S q,t\right)=0}.
\]
$S$ 是 Hamilton 主函数。若 $H$ 不显含时间，
$S=W(q,\alpha)-Et$，得到定态 HJ 方程
\[
H(q,\partial_qW)=E.
\]
在 WKB 中令 $\psi=Ae^{iS/\hbar}$，Schrödinger 方程的
$\hbar\to0$ 领先阶正是 Hamilton--Jacobi 方程，说明经典作用量是量子相位。

\section{小振动的质量加权推导}
稳定平衡附近
\[
T=\frac12\dot{\vb q}^{T}M\dot{\vb q},\qquad
V=\frac12\vb q^TK\vb q.
\]
方程
$M\ddot{\vb q}+K\vb q=0$。因 $M$ 正定，定义
$\vb x=M^{1/2}\vb q$：
\[
\ddot{\vb x}+D\vb x=0,\qquad
D=M^{-1/2}KM^{-1/2}.
\]
$D$ 为实对称矩阵，存在正交本征矢
$D\vb e_\lambda=\omega_\lambda^2\vb e_\lambda$。展开
$\vb x=\sum_\lambda Q_\lambda\vb e_\lambda$：
\[
L=\frac12\sum_\lambda
(\dot Q_\lambda^2-\omega_\lambda^2Q_\lambda^2).
\]
因此正规模是同时对角化动能和势能的坐标。原坐标本征矢满足质量内积正交：
\[
\vb a_\lambda^TM\vb a_{\lambda'}=\delta_{\lambda\lambda'}.
\]

\section{刚体 Euler 方程}
惯性系与随体系导数满足
\[
\left(\dv{\vb A}{t}\right)_{\rm space}
=\left(\dv{\vb A}{t}\right)_{\rm body}
+\vb\omega\times\vb A.
\]
取 $\vb A=\vb L$，Newton 转动定律
$(\dd\vb L/\dd t)_{\rm space}=\vb N$：
\[
\left(\dv{\vb L}{t}\right)_{\rm body}
+\vb\omega\times\vb L=\vb N.
\]
主轴系 $\vb L=(I_1\omega_1,I_2\omega_2,I_3\omega_3)$：
\[
I_1\dot\omega_1+(I_3-I_2)\omega_2\omega_3=N_1
\]
及循环置换。无力矩时绕最大或最小主惯量轴稳定，绕中间轴不稳定。

\begin{selftestbox}
\begin{enumerate}
\item 推导惯性项恒等式时哪两项抵消？
\item 变分端点若不固定会出现什么自然边界条件？
\item Noether 变换使 $L$ 改变全导数为何仍有守恒量？
\item Legendre 变换何时不可逆？
\item Liouville 定理为何不等于“轨迹不混合”？
\item 从正则一形式推导第二类生成函数关系。
\item HJ 方程怎样连接 WKB？
\item 质量矩阵内积下正规模为何正交？
\item Euler 方程中的叉积项从哪里来？
\item 中间轴不稳定的线性化思路是什么？
\end{enumerate}
\end{selftestbox}
